\paragraph{Lee--Yang 三次闭包、Latt\`es 六次正交、\texorpdfstring{$\chi^2$}{chi2} 缺口与 Jacobian 半单刚性}
在 Lee--Yang 分支不变量的共同三次数域、bin-fold 碰撞缺口的 Fibonacci 共振强制，以及 \(S_4\)--Hurwitz Jacobian 的可见等型块分解已经确立之后，还可进一步抽取出下列一组彼此兼容的刚性后果。

\begin{theorem}[Lee--Yang 分支几何的单参数有理统一化]\label{thm:xi-time-part65b-leyang-branch-b-birational-uniformization}
\leanverified{paper\_killo\_leyang\_b\_square\_cubic\_and\_rational\_intertwine}
沿用定理 \ref{thm:xi-terminal-zm-b2-kappa2-rational-interconstruction}、定理 \ref{thm:xi-terminal-zm-stokes-t-elliptic-branch-coordinate} 与定理 \ref{thm:xi-terminal-zm-leyang-elliptic-structure} 的记号。令 \((\lambda_*,y_*)\) 为非平凡 Lee--Yang 分歧点，并记
\[
u:=\kappa_*^2,
\qquad
b:=B^2.
\]
则 \(u\)、\(\lambda_*\) 与 \(y_*\) 都可写成 \(b\) 的有理函数；具体地，
\[
u=\frac{2294-8463b-918b^2}{14388},
\]
\[
\lambda_*
=
\frac{918b^2+8463b+4900}{2(918b^2+8463b+7298)},
\]
\[
y_*
=
-\frac{3597\,b\,(918b^2+8463b+4900)^2}
{(918b^2+8463b-2294)(918b^2+8463b+7298)^2}.
\]
因此
\[
\QQ(b)=\QQ(u)=\QQ(\lambda_*)=\QQ(y_*).
\]
换言之，Lee--Yang 分支点的局部几何不变量 \(u,\lambda_*,y_*,b\) 并非只是在同一三次数域中同居，而可由单一参数 \(b\) 作显式有理统一化。
\end{theorem}
\begin{proof}
由定理 \ref{thm:xi-terminal-zm-b2-kappa2-rational-interconstruction}，
\[
b
=
-2\cdot
\frac{24708u^2-28628u+7733}{51084u^2-44412u+8735}.
\]
同一定理并已给出 \(\QQ(b)=\QQ(u)\)，故 \(u\in \QQ(b)\) 必可唯一写成
\[
u=x_0+x_1b+x_2b^2
\qquad (x_0,x_1,x_2\in \QQ).
\]
把这一待定表达代回上式，并结合 \(b\) 的三次关系
\[
162b^3+1593b^2+1998b+1184=0
\]
在 \(\QQ(b)\) 中约化，即唯一解得
\[
u=\frac{2294-8463b-918b^2}{14388}.
\]

再由定理 \ref{thm:xi-terminal-zm-kappa-square-cubic-field-s3}，
\[
\lambda_*=\frac{3(2u-1)}{4(3u-2)}.
\]
把前式代入并化简，便得到所示 \(\lambda_*(b)\) 的表达。

另一方面，定理 \ref{thm:xi-terminal-zm-stokes-t-elliptic-branch-coordinate} 给出
\[
4\lambda_*y_*=4\lambda_*^3-3\lambda_*^2-2\lambda_*+1.
\]
再把 \(\lambda_*(b)\) 代入并整理，即得所示 \(y_*(b)\) 的表达。

由定理 \ref{thm:xi-terminal-zm-b2-kappa2-rational-interconstruction} 已知
\[
\QQ(b)=\QQ(u)=\QQ(\lambda_*).
\]
又由于 \(y_*\) 是定理 \ref{thm:xi-time-part9g-leyang-cubic-discriminant} 中不可约三次
\[
P_{\mathrm{LY}}(y)=256y^3+411y^2+165y+32
\]
的根，故 \([\QQ(y_*):\QQ]=3\)。而上式已证明 \(y_*\in \QQ(b)\)，于是只能有
\[
\QQ(y_*)=\QQ(b).
\]
遂得全部域相等式。证毕。
\end{proof}

\begin{theorem}[Lee--Yang 分支不变量链的共同二次分辨子域与统一 \texorpdfstring{$S_3$}{S3} 闭包]\label{thm:xi-time-part65b-leyang-common-s3-closure}
沿用定理 \ref{thm:xi-time-part65b-leyang-branch-b-birational-uniformization} 的记号，记
\[
u:=\kappa_*^2,
\qquad
b:=B^2,
\qquad
K:=\QQ(u)=\QQ(b)=\QQ(\lambda_*)=\QQ(y_*).
\]
设 \(u\) 与 \(b\) 的三次最小多项式分别为
\[
324u^3-540u^2+333u-74,
\qquad
162b^3+1593b^2+1998b+1184,
\]
其判别式分别记为
\[
\Delta_u=-2^6\cdot 3^9\cdot 37,
\qquad
\Delta_b=-2^2\cdot 3^9\cdot 11^2\cdot 37\cdot 109^2.
\]
若 \(L_{\mathrm{LY}}\) 记为三次数域 \(K/\QQ\) 的伽罗瓦闭包，则有
\[
\mathrm{sf}(\Delta_u)=\mathrm{sf}(\Delta_b)=-111,
\]
并且
\[
\Gal(L_{\mathrm{LY}}/\QQ)\cong S_3,
\qquad
L_{\mathrm{LY}}^{A_3}=\QQ(\sqrt{-111}).
\]
等价地，\(u\)、\(b\)、\(\lambda_*\) 与 \(y_*\) 都落在同一非伽罗瓦三次数域 \(K\) 中，其共同伽罗瓦闭包唯一的非平凡二次正规子域恰为 \(\QQ(\sqrt{-111})\)。
\end{theorem}
\begin{proof}
由定理 \ref{thm:xi-time-part65b-leyang-branch-b-birational-uniformization} 与定理 \ref{thm:xi-terminal-zm-stokes-leyang-shared-artin-representation}，
\[
\QQ(u)=\QQ(b)=\QQ(\lambda_*)=\QQ(y_*)=:K,
\qquad
\Gal(L_{\mathrm{LY}}/\QQ)\cong S_3.
\]
因此 \(u\)、\(b\)、\(\lambda_*\) 与 \(y_*\) 都生成同一非伽罗瓦三次数域 \(K\)，其分裂域必同为 \(L_{\mathrm{LY}}\)。

另一方面，
\[
\Delta_u
=
-2^6\cdot 3^9\cdot 37
=
(2^3\cdot 3^4)^2\cdot (-111),
\]
\[
\Delta_b
=
-2^2\cdot 3^9\cdot 11^2\cdot 37\cdot 109^2
=
(2\cdot 3^4\cdot 11\cdot 109)^2\cdot (-111),
\]
故
\[
\mathrm{sf}(\Delta_u)=\mathrm{sf}(\Delta_b)=-111.
\]
对不可约三次域，其伽罗瓦闭包的唯一非平凡二次正规子域由判别式的平方自由核决定，因此
\[
L_{\mathrm{LY}}^{A_3}=\QQ(\sqrt{-111}).
\]
证毕。
\end{proof}

\begin{proposition}[Lee--Yang 三次数域的两组显式整生成元]\label{prop:xi-time-part65b-leyang-u-b-integral-generators}
沿用定理 \ref{thm:xi-time-part65b-leyang-common-s3-closure} 的记号，定义
\[
x:=18u,
\qquad
y:=18b.
\]
则 \(x\) 与 \(y\) 都是代数整数，并分别满足首一整系数三次方程
\[
x^3-30x^2+333x-1332=0,
\]
\[
y^3+177y^2+3996y+42624=0.
\]
特别地，
\[
\QQ(x)=\QQ(u)=K,
\qquad
\QQ(y)=\QQ(b)=K.
\]
\end{proposition}
\begin{proof}
由定理 \ref{thm:xi-time-part65b-leyang-common-s3-closure}，
\[
324u^3-540u^2+333u-74=0.
\]
代入 \(u=x/18\) 并清分母，得到
\[
x^3-30x^2+333x-1332=0.
\]
同理，由
\[
162b^3+1593b^2+1998b+1184=0
\]
代入 \(b=y/18\) 并清分母，得到
\[
y^3+177y^2+3996y+42624=0.
\]
两式皆为首一整系数多项式，故 \(x\) 与 \(y\) 都是代数整数。
又因 \(18\in \QQ^\times\)，缩放不改变所生成的 \(\QQ\)-子域，因此
\[
\QQ(x)=\QQ(u)=K,
\qquad
\QQ(y)=\QQ(b)=K.
\]
证毕。
\end{proof}

\begin{corollary}[共同 \texorpdfstring{$S_3$}{S3} 闭包与 Latt\`es 六次域的线性互素性]\label{cor:xi-time-part65b-leyang-lattes-linear-disjointness}
设
\[
C(X)=X^6-8X^5+24X^4-36X^3+34X^2-20X+4
\]
的分裂域为 \(L_4\)，并满足
\[
\Disc(C)=-2^{19}\cdot 37^3,
\qquad
\Gal(L_4/\QQ)\cong S_4\times C_2.
\]
则定理 \ref{thm:xi-time-part65b-leyang-common-s3-closure} 中的共同闭包 \(L_{\mathrm{LY}}\) 与 \(L_4\) 线性互素：
\[
L_{\mathrm{LY}}\cap L_4=\QQ.
\]
因而
\[
\Gal(L_{\mathrm{LY}}L_4/\QQ)
\cong
S_3\times (S_4\times C_2).
\]
\end{corollary}
\begin{proof}
由定理 \ref{thm:xi-time-part65b-leyang-common-s3-closure}，
\[
L_{\mathrm{LY}}^{A_3}=\QQ(\sqrt{-111}),
\]
这是 \(L_{\mathrm{LY}}/\QQ\) 的唯一非平凡二次正规子域。于是交域
\[
F:=L_{\mathrm{LY}}\cap L_4
\]
作为 \(L_{\mathrm{LY}}/\QQ\) 的伽罗瓦子扩张，只能是 \(\QQ\) 或 \(\QQ(\sqrt{-111})\)。

若 \(F=\QQ(\sqrt{-111})\)，则素数 \(3\) 在 \(F/\QQ\) 中分歧；然而由
\[
\Disc(C)=-2^{19}\cdot 37^3
\]
可知 \(L_4/\QQ\) 只在 \(\{2,37\}\) 处分歧，故 \(3\) 在 \(L_4\) 的任一子扩张中都不分歧，矛盾。故只能有
\[
L_{\mathrm{LY}}\cap L_4=\QQ.
\]
最后，两边皆为 \(\QQ\) 上的伽罗瓦扩张，故复合域的伽罗瓦群为直积：
\[
\Gal(L_{\mathrm{LY}}L_4/\QQ)
\cong
\Gal(L_{\mathrm{LY}}/\QQ)\times \Gal(L_4/\QQ)
\cong
S_3\times (S_4\times C_2).
\]
证毕。
\end{proof}

\begin{theorem}[临界共振常数强制 \texorpdfstring{$\chi^2$}{chi2} 碰撞缺口具有正的极限下界]\label{thm:xi-time-part65b-cphi-chi2-gap-floor}
沿用定义 \ref{def:fold-normalized-amplitude}、命题 \ref{prop:fold-bin-chi2-col} 与定理 \ref{thm:fold-critical-resonance-constant} 的记号。记
\[
a_m(k):=\frac{|\widehat d_m(k)|}{2^m},
\qquad
C_\varphi:=\lim_{m\to\infty}a_m(F_m)=\lim_{m\to\infty}a_m(F_{m+1})>0.
\]
若 \(p_m^{\mathrm{bin}}\) 为 bin-fold 推前分布，\(u_m\) 为 \(X_m\) 上的均匀分布，则
\[
\liminf_{m\to\infty}\chi^2\!\bigl(p_m^{\mathrm{bin}}\|u_m\bigr)\ge 2C_\varphi^2.
\]
等价地，
\[
\liminf_{m\to\infty}
F_{m+2}\Bigl(\mathsf{Col}_m-\frac{1}{F_{m+2}}\Bigr)
\ge
2C_\varphi^2.
\]
\end{theorem}
\begin{proof}
由命题 \ref{prop:fold-bin-chi2-col}，
\[
\chi^2\!\bigl(p_m^{\mathrm{bin}}\|u_m\bigr)
=
|X_m|\,\mathsf{Col}_m-1
=
F_{m+2}\Bigl(\mathsf{Col}_m-\frac{1}{F_{m+2}}\Bigr).
\]
另一方面，结论 \ref{con:xi-time-part9m-cphi-forces-collision-gap-and-peak-bias} 已给出
\[
\liminf_{m\to\infty}
F_{m+2}\Bigl(\mathsf{Col}_m-\frac{1}{F_{m+2}}\Bigr)
\ge
2C_\varphi^2.
\]
将二者合并即得
\[
\liminf_{m\to\infty}\chi^2\!\bigl(p_m^{\mathrm{bin}}\|u_m\bigr)\ge 2C_\varphi^2.
\]
第二式只是上述恒等式的等价改写。证毕。
\end{proof}

\begin{theorem}[可见等型块代数的显式半单嵌入与 N\'eron--Severi 秩下界]\label{thm:xi-time-part65b-jacobian-visible-semisimple-ns-floor}
沿用定理 \ref{thm:killo-s4-genus49-jacobian-computable-isogeny} 的 \(\QQ\)-同源分解
\[
J(X)\sim_{\QQ}J(H)\times J(Z)^2\times J(C_F)^3\times P^3.
\]
则 \(\operatorname{End}_{\QQ}^0(J(X))\) 含有显式半单 \(\QQ\)-子代数
\[
\mathcal A_X:=\QQ\times M_2(\QQ)\times M_3(\QQ)\times M_3(\QQ)
\hookrightarrow
\operatorname{End}_{\QQ}^0(J(X)),
\]
从而
\[
\dim_{\QQ}\operatorname{End}_{\QQ}^0(J(X))\ge 1+4+9+9=23.
\]
进一步，在右端乘积上取任意乘积极化，并把相应 Rosati 对合记为 \(\dagger_{\mathrm{vis}}\)。则
\[
\dim_{\QQ}(\mathcal A_X)^{\dagger_{\mathrm{vis}}}
=
1+\frac{2\cdot 3}{2}+\frac{3\cdot 4}{2}+\frac{3\cdot 4}{2}
=
16.
\]
特别地，
\[
\operatorname{rank}\operatorname{NS}(J(X)_{\overline{\QQ}})\ge 16.
\]
\end{theorem}
\begin{proof}
定理 \ref{thm:xi-time-part9n1b-jacobian-projector-idempotent-algebra} 已经给出显式嵌入
\[
\mathcal A_X
\hookrightarrow
\operatorname{End}_{\QQ}^0(J(X))
\]
以及
\[
\dim_{\QQ}\operatorname{End}_{\QQ}^0(J(X))\ge 23.
\]
故只需证明 Rosati 不动维数与 N\'eron--Severi 秩下界。

取一个实现上述 \(\QQ\)-同源分解的同源映射，并在乘积
\[
A\times B^2\times C^3\times D^3
\qquad
\bigl(A:=J(H),\ B:=J(Z),\ C:=J(C_F),\ D:=P\bigr)
\]
上选取乘积极化 \(\lambda_{\mathrm{prod}}\)。在该模型下，\(\mathcal A_X\) 恰由四个可见等型块的重数空间上的标量矩阵作用给出，因此 \(\lambda_{\mathrm{prod}}\) 的 Rosati 对合在 \(\QQ\) 因子上恒等，在 \(M_r(\QQ)\) 因子上与转置等价。于是其不动子空间维数分别为
\[
1,\qquad \frac{2\cdot 3}{2}=3,\qquad \frac{3\cdot 4}{2}=6,\qquad \frac{3\cdot 4}{2}=6,
\]
从而
\[
\dim_{\QQ}(\mathcal A_X)^{\dagger_{\mathrm{vis}}}=1+3+6+6=16.
\]

最后，同源映射的拉回在 N\'eron--Severi 群上于有理化后给出同构，而对任意极化阿贝尔簇 \(A\) 均有标准识别
\[
\operatorname{NS}(A_{\overline{\QQ}})\otimes_{\ZZ}\QQ
\cong
\operatorname{End}^0_{\overline{\QQ}}(A)^{\dagger=1}.
\]
故上述 \(16\) 维 Rosati 不动子空间传回 \(J(X)\) 后仍给出
\[
\operatorname{rank}\operatorname{NS}(J(X)_{\overline{\QQ}})\ge 16.
\]
证毕。
\end{proof}

\noindent
以上五条结果把 Lee--Yang 分支几何的单参数有理统一化、共同 \(S_3\)-闭包、Latt\`es 六次闭包、黄金共振二阶缺口与 \(S_4\)--Hurwitz Jacobian 的可见块分解进一步压缩到同一层级：一方面，\(b\) 不仅与 \(u\)、\(\lambda_*\) 及 \(y_*\) 生成同一三次数域，而且把全部分支点局部几何显式统一化；该共同三次数域的伽罗瓦闭包唯一的非平凡二次正规子域仍被 \(\QQ(\sqrt{-111})\) 完全锁定，并与 \(L_4\) 严格正交；另一方面，临界共振常数把 bin-fold 与均匀基线之间的 \(\chi^2\) 偏离锁定为严格正的极限底噪，而 Jacobian 的可见等型块则单独强制出一个 \(23\) 维有理半单子代数与至少 \(16\) 维的 N\'eron--Severi 秩下界。

\endinput
